\chapter{Qué es, exactamente, un jacobiano secuencial}
\label{cap:jacobiano}

\section{Definición y lectura}
\label{sec:jacobiano-def}

\begin{definicion}[Jacobiano secuencial]
Sea $h$ una función que mapea la trayectoria de un insumo agregado
$\mathbf{X} = (X_0,\dots,X_{T-1})$ a la trayectoria de un producto agregado
$\mathbf{Y} = (Y_0,\dots,Y_{T-1})$, y supongamos que $h$ es diferenciable en el
estado estacionario $\mathbf{X}_\sst$. El \emph{jacobiano secuencial} de $h$ es
la matriz $T\times T$
\begin{equation}
\J_{t,s} \;=\; \frac{\partial Y_t}{\partial X_s}\bigg|_{\mathbf{X}=\mathbf{X}_\sst}.
\end{equation}
\end{definicion}

La interpretación que vuelve útil a este objeto es la de \emph{noticia}
(\emph{news shock}). La columna $s$ de $\J$ es la respuesta de toda la
trayectoria de $Y$ cuando, en la fecha $0$, los agentes se enteran de que $X$
será una unidad más alto en la fecha $s$ y solo en esa fecha. Por eso cada
columna se lee como una función impulso-respuesta, y por eso la matriz completa
contiene las respuestas a \emph{todas} las trayectorias posibles de $X$: una
trayectoria arbitraria $d\mathbf{X}$ es una combinación lineal de noticias, y
\begin{equation}
d\mathbf{Y} = \J\, d\mathbf{X}
\end{equation}
es un simple producto matriz--vector. Calcular una función impulso-respuesta
adicional, una vez que se tiene $\J$, es gratis.

\section{Un ejemplo que se puede hacer a mano}

Antes de ir al caso heterogéneo, conviene ver un jacobiano que se puede derivar
con lápiz. Tomemos un consumidor representativo con horizonte infinito,
utilidad $\sum \beta^t u(c_t)$ con $u(c)=\log c$, que recibe un ingreso
exógeno $y_t$ y puede ahorrar a una tasa constante $r$ con $\beta(1+r)=1$. Su
restricción presupuestaria intertemporal es
\begin{equation}
\sum_{t\ge 0} \frac{c_t}{(1+r)^t} = (1+r)a_{-1} + \sum_{t\ge 0}\frac{y_t}{(1+r)^t}.
\end{equation}
Con utilidad logarítmica y $\beta(1+r)=1$ el consumo es constante e igual a la
anualidad de la riqueza:
\begin{equation}
c_t = \frac{r}{1+r}\Big[(1+r)a_{-1} + \sum_{s\ge0}\frac{y_s}{(1+r)^s}\Big].
\end{equation}
Derivando,
\begin{equation}
\J^{C,y}_{t,s} = \frac{\partial c_t}{\partial y_s} = \frac{r}{1+r}\,\frac{1}{(1+r)^s}.
\label{eq:jac-ri}
\end{equation}

Este jacobiano tiene todas sus filas iguales: el agente suaviza perfectamente,
así que la fecha en que llega el ingreso no afecta en qué momento lo consume,
solo cuánto vale en valor presente. La propensión marginal a consumir del
ingreso contemporáneo es $r/(1+r)$, que con $r=1\%$ trimestral es $0.0099$: un
céntimo por cada sol. Es el resultado clásico de la hipótesis de ingreso
permanente, y es \emph{exactamente} lo que los datos rechazan.

\begin{idea}
El jacobiano $\J^{C,y}$ de un modelo de agentes heterogéneos es muy distinto del
de \eqref{eq:jac-ri}: su diagonal es grande (propensiones marginales a consumir
de 15 a 30 céntimos por sol, no uno), y decae rápido fuera de ella. Esa
diferencia --- y no ninguna otra --- es la que hace que la política fiscal y
monetaria funcionen distinto en un modelo HANK que en uno representativo
(Auclert, Rognlie y Straub, 2018). Comparar jacobianos es la forma más limpia de
comparar modelos.
\end{idea}

\section{El jacobiano de un modelo de mercados incompletos}

La figura \ref{fig:columnas} muestra cuatro columnas del jacobiano del ahorro y
del consumo agregados respecto de la tasa de interés, en el modelo de
Krusell--Smith calibrado trimestralmente ($\alpha=0.36$, $\delta=0.025$,
$K/Y=10$, elasticidad de sustitución intertemporal $0.5$, proceso de ingreso
AR(1) con $\rho=0.966$). Todo lo que aparece en esta y en las demás figuras sale
del código que acompaña al manual.

\begin{figure}[htbp]
\centering
\includegraphics[width=\textwidth]{fig01_columnas_jacobiano.pdf}
\caption{Columnas del jacobiano del bloque de hogares respecto de la tasa de
interés, en el modelo de Krusell--Smith. Cada curva es la respuesta de la
variable agregada a la \emph{noticia}, recibida en la fecha $0$, de que la tasa
será un punto más alta en la fecha $s$. El patrón de ``carpa'' --- acumulación
previa al choque, pico en $s$, desacumulación posterior --- es la firma
característica de estos modelos.}
\label{fig:columnas}
\end{figure}

Vale la pena leer la figura con calma, porque contiene casi toda la economía del
modelo.

\paragraph{La columna $s=0$} Cuando la tasa sube hoy, sin aviso, los hogares
reciben un retorno mayor sobre los activos que ya tenían. Es un efecto riqueza
puro: ahorran parte de ese ingreso extra, y en los períodos siguientes lo van
desacumulando lentamente. La curva arranca en su máximo y decae.

\paragraph{Las columnas $s>0$} Aquí aparece la sustitución intertemporal. Si el
hogar sabe que la tasa será alta dentro de 50 trimestres, le conviene llegar a
esa fecha con más activos. Entonces ahorra desde hoy, acumulando
progresivamente, llega al pico justo en $s$, y después desacumula. De ahí la
``carpa''. La parte ascendente es anticipación; la descendente, el mismo efecto
riqueza de antes.

\paragraph{El consumo} El panel derecho muestra el espejo. Ante una noticia de
tasa alta futura, el consumo \emph{cae} hoy (el agente está ahorrando para
llegar con más activos) y salta en la fecha del choque. La caída inicial es
pequeña en relación con el salto: los hogares restringidos no pueden sustituir
intertemporalmente, y en esta calibración un 2.5\% de ellos están contra la
restricción de endeudamiento en todo momento.

\section{La estructura que hace posible el algoritmo rápido}

La figura \ref{fig:estructura} muestra la misma matriz de dos maneras.

\begin{figure}[htbp]
\centering
\includegraphics[width=\textwidth]{fig02_estructura_jacobiano.pdf}
\caption{Izquierda: mapa de calor de $\J^{A,r}$. Derecha: las columnas
$s=40,70,100,130$ realineadas en $t-s$. Lejos del origen, las columnas son
esencialmente traslados unas de otras: el jacobiano es asintóticamente
invariante en el tiempo.}
\label{fig:estructura}
\end{figure}

El panel derecho contiene la observación crucial. Si uno toma la columna $s=70$ y
la desplaza 30 períodos hacia la izquierda, queda casi exactamente encima de la
columna $s=100$. Formalmente,
\begin{equation}
\J_{t,s} \;\approx\; \J_{t-1,s-1} \qquad \text{para } t,s \text{ grandes}.
\label{eq:asint}
\end{equation}

¿Por qué? Dos razones independientes.

\begin{enumerate}[leftmargin=1.4em]
\item \textbf{Hacia atrás:} el comportamiento de un agente ante una noticia
depende solo de la \emph{distancia} a esa noticia, no de la fecha del
calendario. Cómo reacciona hoy a algo que pasará en 50 trimestres es lo mismo
que cómo reaccionará mañana a algo que pasará 50 trimestres después de mañana.
Esto es exacto, no aproximado, y es el lema \ref{lema:politicas} del capítulo
\ref{cap:fakenews}.
\item \textbf{Hacia adelante:} el efecto de un cambio en la distribución se
disipa. Si hoy la distribución se perturba un poco, dentro de 200 trimestres
casi no queda rastro, porque el proceso de ingreso es ergódico y la mezcla
elimina la información inicial. La figura \ref{fig:expectativa} cuantifica esa
disipación.
\end{enumerate}

La relación \eqref{eq:asint} no es solo una curiosidad estética. Significa que
la matriz $\J$, que tiene $T^2 = 90\,000$ entradas, está determinada por un
número mucho menor de objetos independientes. El algoritmo fake news es,
precisamente, el procedimiento que explota esa redundancia para calcular las
$90\,000$ entradas al costo de calcular unas $2T$.
